\section{Conclusion}
\label{submission:sec:conclusion}

In this paper, we have adopted the critical perspective of {\em The Methodologist} to systematically deconstruct the out-of-distribution (OOD) task rejection modules inside dynamic model merging frameworks. We exposed a hidden yet critical vulnerability: unregularized coordinate-space density models (e.g., GMMs) fit on low-resource calibration splits ($N \le 64$) suffer from severe variance collapse. This collapse makes their log-likelihood boundaries highly unstable, causing catastrophic false positive rate spikes under realistic, representation-level covariate shifts. 

To resolve this methodological issue, we introduced {\bf SRC-DE} (Shrinkage-Regularized Coordinate Density Estimation), which applies analytical, parameter-free Ledoit-Wolf-style covariance shrinkage as a post-fit regularization step immediately following EM convergence. Our rigorous empirical evaluation across high-conflict vision task suites under varying levels of representation noise and calibration sample sizes validated SRC-DE's ability to suppress variance under small sample splits. Under moderate calibration splits ($N=32$), unregularized diagonal GMMs exhibit severe estimation instability with a massive standard deviation of $\pm 0.0863$ (AUC of $0.6932 \pm 0.0863$). By analytically computing the optimal shrinkage intensity, our proposed SRC-DE successfully regularizes coordinate covariances, stabilizing the log-likelihood boundaries and achieving an AUC of $0.7346 \pm 0.0632$. Under extreme data scarcity ($N=8$), SRC-DE provides a critical stability guarantee, yielding a standard deviation of $\pm 0.0270$ (AUC of $0.7781 \pm 0.0270$), compared to $\pm 0.0292$ (AUC of $0.7683 \pm 0.0292$) for the unregularized GMM. Furthermore, we identified and resolved a critical, silent bug in standard python GMM libraries regarding manual covariance modifications.

\subsection{Generalization to Other Modalities and Scale}
While our empirical validation focuses on edge-friendly vision tasks and a lightweight ViT backbone, the mathematical formulation of similarity coordinates and responsibility-weighted covariance shrinkage is modal-agnostic and scales seamlessly. In natural language processing (NLP) task registries, serving networks route client prompts across task-specific parameter-efficient adapters (such as LoRA experts) stacked on heavy, frozen language backbones (such as LLaMA or Mistral). In these architectures, early-layer hidden representations (such as prompt CLS token activations or average-pooled token states) can be projected onto task centroids in exactly the same manner to derive task similarity coordinates. A comprehensive discussion on generalization to convolutional backbones (e.g., ConvNeXt) and heavier vision transformers (e.g., ViT-Large) is provided in Appendix~\ref{sec:backbone_generalization}.

As task registries scale to larger, multi-tenant serving networks supporting dozens of experts ($K \ge 50$), the coordinate-space dimension $K$ expands, while the calibration budget $N$ per task remains severely limited due to data privacy or device constraints. In these high-dimensional, multi-modal registries, unregularized density boundaries will overfit exponentially fast due to the curse of dimensionality. Applying SRC-DE's responsibility-weighted Ledoit-Wolf shrinkage in high-dimensional coordinate spaces represents a mathematically mandatory path to enable expressive, multi-modal routing while guaranteeing absolute statistical and numerical stability.

\subsection{Future Research Directions}
To guide the next generation of dynamic serving and statistical routing research, we identify three highly promising, theoretically rich directions for future investigation:
\begin{itemize}
    \item {\bf Dynamic Noise Estimation (Oracle Mitigation):} While our highly competitive Noise-Adapted SRC-DE variant (yielding $0.8137$ AUC under overlap) relies on an oracle test-time representation noise variance $\sigma^2$, real-world serving environments require a dynamic, data-driven approach. Practitioners can estimate this noise on the fly by tracking the running variance of similarity coordinates from recently routed, high-confidence queries relative to their nearest task centroids, providing an adaptive, hyperparameter-free noise adjustment without prior oracle knowledge.
    \item {\bf Delta-Method for EM Responsibility Sampling Variance:} As detailed in Appendix~A.1, treating GMM posterior responsibilities $\gamma_{s, m}$ as fixed constants underestimates the true variance of coordinate variance estimators, particularly under EM splitting instability ($N \le 16$). Integrating delta-method approximations or local bootstrap estimators directly into the M-step of EM could capture this responsibility sampling variance, mathematically refining the analytical shrinkage intensity $\alpha_{\text{opt}}$ to achieve absolute statistical tightness.
    \item {\bf Hybrid Bounded-Support Density Estimators:} Because cosine similarity coordinates reside strictly on the bounded support $[-1, 1]$, standard Gaussian components inevitably suffer from minor probability mass leakage beyond the boundaries. Although Beta mixtures and truncated Gaussians formally address this, their optimization is highly unstable under extreme data scarcity. A highly promising compromise is a hybrid bounded estimator, where standard, closed-form Gaussian covariance components are analytically shrunk toward truncated Gaussian or Beta-mixture priors, merging the low-resource optimization stability of unconstrained Gaussians with the theoretical rigor of bounded supports.
\end{itemize}

By exposing the clean sandbox confounder and the low-resource variance collapse, this work emphasizes the importance of methodological skepticism in model merging research. We hope our findings encourage the community to adopt more rigorous evaluation protocols, and that SRC-DE serves as a mathematically sound and practical foundation for robust, sample-efficient dynamic model serving on the edge.
